\chapter{Refinement criteria and conservative transfer}\label{ch:amr-criteria-transfer}
\chapterstatus{outline}
Thesis Chapter~2 describes refinement criteria but not this package's
conservative block-tree transfer~\cite{remmerswaal2023}. The chapter will map
indicators and derive cell-integral, divergence, and phase-transfer properties.

\section{Refinement criteria}
The available interface, curvature, vorticity, geometry, and interpolation
error criteria will be compared with the thesis's tolerances.
\begin{deviation}{Interface refinement}
The package uses a target level and buffer; the thesis's interface criterion uses a phase-fraction threshold.
\end{deviation}
\begin{deviation}{Instability refinement}
The thesis's instability-wavelength refinement criterion is not implemented.
\end{deviation}
\begin{implnote}
\apitype{InterfaceRefinementIndicator}, \apitype{CurvatureRefinementIndicator},
\apitype{VorticityRefinementIndicator}, \apitype{InterpolationErrorIndicator},
\apitype{GeometryRefinementIndicator}.
\end{implnote}

\section{Cell, face, and phase transfer}
Cell restriction and prolongation, face transfer, and PLIC/phase transfer will
be related to their conservation requirements.
\begin{implnote}
\apifun{restrict!}, \apifun{prolong!}, \apitype{AMRWorkspace}.
\end{implnote}

\section{Divergence preservation and regridding}
The local Cartesian-divergence correction is applied on refinement; coarsening
uses area-averaged face velocities. In a two-velocity flow simulation the
phase-aperture-weighted mixture is projected after either topology change, as
the independently transferred phase fields need not preserve its divergence
(\cref{eq:vof:frozen-flux}). Proofs of the transfer invariants belong here.
\begin{implnote}
\apitype{MeshAdaptationOptions}, \apifun{adapt!}, \apifun{divergence!}.
\end{implnote}
